\chapter{Glossary}
\label{ch:glossary}

\begin{description}
\item[AGM] \textbf{Arithmetic-Geometric Mean.} The common limit of the sequences $a_{n+1} = \frac{a_n+b_n}{2}$ and $b_{n+1} = \sqrt{a_n b_n}$, introduced by Gauss to relate elliptic integrals and $\pi$.

\item[Binary Quadratic Form] A polynomial $ax^2 + bxy + cy^2$ with integer coefficients. Gauss developed the composition theory for these forms in the \textit{Disquisitiones Arithmeticae}.

\item[Chinese Remainder Theorem] A theorem stating that a system of congruences with pairwise coprime moduli has a unique solution modulo the product of the moduli. Gauss gave the first complete proof.

\item[Class Number] The number of equivalence classes of binary quadratic forms of a given discriminant, or equivalently the order of the ideal class group of a quadratic field.

\item[Congruence] The relation $a \equiv b \pmod{n}$ meaning $n \mid (a-b)$. Gauss introduced this notation and systematized its theory.

\item[Dedekind Eta Function] A modular form defined by $\eta(\tau) = e^{\pi i \tau/12} \prod_{n=1}^{\infty}(1-e^{2\pi i n\tau})$, closely related to theta functions.

\item[Discriminant] For a quadratic form $ax^2+bxy+cy^2$, the value $\Delta = b^2-4ac$. For a polynomial, the product of squared differences of roots.

\item[Dirichlet Character] A completely multiplicative periodic function $\chi: \Z \to \C$ with $\chi(n) = 0$ when $\gcd(n,m) > 1$, used to define Dirichlet $L$-functions.

\item[Error Function] The function $\operatorname{erf}(x) = \frac{2}{\sqrt{\pi}}\int_0^x e^{-t^2}\,dt$, fundamental in probability and the study of Gaussian distributions.

\item[Fermat Prime] A prime of the form $2^{2^k}+1$. Gauss proved that a regular $n$-gon is constructible iff $n$ is a product of distinct Fermat primes and a power of 2.

\item[Gaussian Elimination] A method for solving systems of linear equations by row reduction. Gauss used this extensively in geodesy and celestial mechanics.

\item[Gaussian Integral] The integral $\int_{-\infty}^{\infty} e^{-x^2}\,dx = \sqrt{\pi}$, central to probability theory and mathematical physics.

\item[Gauss Sum] An exponential sum $g(a,p) = \sum_{x=0}^{p-1} e^{2\pi i ax^2/p}$ or its generalization $\tau(\chi) = \sum_{a=0}^{n-1} \chi(a) e^{2\pi i a/n}$ for Dirichlet characters.

\item[Jacobi Theta Function] One of the four functions $\vartheta_1, \vartheta_2, \vartheta_3, \vartheta_4$ of two variables, satisfying modular transformation laws and product identities.

\item[Jacobian] The determinant of the matrix of first partial derivatives, measuring the scaling factor under a multivariate transformation.

\item[Jacobi Symbol] A generalization of the Legendre symbol to odd positive denominators, defined multiplicatively: $\left(\frac{a}{n}\right) = \prod \left(\frac{a}{p_i}\right)^{e_i}$.

\item[Kronecker Symbol] A further generalization of the Jacobi symbol to all integer denominators, used in the theory of quadratic fields.

\item[Legendre Symbol] The symbol $\left(\frac{a}{p}\right)$ indicating whether $a$ is a quadratic residue modulo an odd prime $p$.

\item[Least Squares] A method for solving overdetermined systems by minimizing the sum of squared residuals, $\min_x \|Ax - b\|^2$. Gauss applied this to orbit determination.

\item[Linear Congruence] An equation of the form $ax \equiv b \pmod{n}$, solvable iff $\gcd(a,n) \mid b$.

\item[Modular Arithmetic] Arithmetic in the ring $\Z/n\Z$ of residue classes modulo $n$, where operations are performed on representatives and results reduced modulo $n$.

\item[Modular Form] A holomorphic function on the upper half-plane satisfying a transformation law under $SL(2,\Z)$ and a growth condition at cusps.

\item[Normal Distribution] The Gaussian probability density $f(x) = \frac{1}{\sigma\sqrt{2\pi}} e^{-(x-\mu)^2/(2\sigma^2)}$, central to statistics and error theory.

\item[Orbital Elements] The six parameters $(a, e, i, \Omega, \omega, M_0)$ that uniquely determine a Keplerian orbit.

\item[Primitive Root] A generator of the multiplicative group $(\Z/p\Z)^*$ for prime $p$. Gauss proved their existence for all prime moduli.

\item[Quadratic Reciprocity] The theorem relating $\left(\frac{p}{q}\right)$ and $\left(\frac{q}{p}\right)$ for distinct odd primes: $\left(\frac{p}{q}\right)\left(\frac{q}{p}\right) = (-1)^{(p-1)(q-1)/4}$.

\item[Quadratic Residue] An integer $a$ such that $x^2 \equiv a \pmod{p}$ has a solution for an odd prime $p$.

\item[Residue Class] An equivalence class $[a]_n = \{b \in \Z : b \equiv a \pmod{n}\}$ in the ring $\Z/n\Z$.

\item[Theorema Egregium] Gauss's ``Remarkable Theorem'' stating that Gaussian curvature is an intrinsic property of a surface, invariant under isometric deformations.

\item[Theta Function] A special function of several complex variables, including Jacobi theta functions and Siegel theta series, central to modular form theory.

\item[Weierstrass $\wp$-Function] An elliptic function defined by $\wp(z) = \frac{1}{z^2} + \sum_{\omega \in \Lambda\setminus\{0\}} \left(\frac{1}{(z-\omega)^2} - \frac{1}{\omega^2}\right)$, satisfying $(\wp')^2 = 4\wp^3 - g_2\wp - g_3$.

\end{description}
